\section{Sampling, seeds and complexity}\label{sec:complexity}

This section is informal. Nothing in it is formalized beyond the bias bounds it invokes.

\paragraph{Computational model.}
Integers are given in binary. The modulus $q$ is given only as an integer: neither its factorization,
nor $\varphi(q)$, nor any generator of $\Gq$ is needed. The only use of $q$ is a coprimality test.

\paragraph{The sample space is $q$-independent.}
For fixed $(a,y)$ the weights $w_{a,y}(n)$ do not depend on $q$; $q$ enters only through the filter
$\gcd(n,q)=1$. Let $W=\{n\in[xe^{-2a},xe^{2a}]:\Lambda(n)>0\}$. A sieve up to $M=\lfloor e^{y+2a}\rfloor+1$
lists $W$ and the weights in $O(M\log\log M)$ arithmetic operations. Each weight needs one evaluation of
$k_a$, a smooth function of one real variable, computed once per point to any desired precision. For the
certificate of Corollary~\ref{cor:cert}, $M=166254$ and $\abs W$ is about $1.5\cdot10^4$.

\paragraph{Exact sampling and seed length.}
To output a sample with $s$ uniform random bits, round the normalized weights to multiples of $2^{-s}$ and
use inverse-CDF lookup in the precomputed table. If the rounding error in total variation is at most $\delta$,
the bias increases by at most $2\delta$. Taking $s=\lceil\log_2\abs W\rceil+\lceil\log_2(1/\delta)\rceil$ suffices.
With support $X\asymp\log^2q/\varepsilon^2$ (Corollary~\ref{cor:asymp}) this gives an $(\varepsilon+2\delta)$-biased
generator for $\Gq$ with seed length
\[
  s=2\log_2\log q+2\log_2(1/\varepsilon)+\log_2(1/\delta)+O(1)
\]
(the $O(1)$ absorbs the leading constant and the prime-power density). Each sample then costs one table lookup and one
gcd with $q$. The lower bound for any $\varepsilon$-biased set in a group of order $\varphi(q)$ is
$\log_2\log q+2\log_2(1/\varepsilon)-O(\log\log(1/\varepsilon))$ \cite{alonroichman1994,naornaor1993}. The
factor $2$ in front of $\log_2\log q$ is the $(\log q)^2$ barrier of Section~\ref{sec:limits}.

\paragraph{Non-units and $q$-obliviousness.}
Sampling from $W$ without the coprimality filter gives a measure that does not depend on $q$ at all. The
non-units of $W$ carry relative mass at most $U_q(a,y)/\mathrm{Den}$, which Theorem~\ref{thm:closed} bounds by
$U_B/\mathrm{den}_B$. For $q\le2^{64}$ with the parameters of Corollary~\ref{cor:cert} this is at most
$37.15/555.8<0.067$. So a single, $q$-oblivious table serves all moduli $q\le2^{64}$ simultaneously, at the
cost of a small failure probability.

\paragraph{Expander walks.}
By Corollary~\ref{cor:cert}, the random walk on $\Gq$ that multiplies by an independent sample at each step has
second eigenvalue at most $\frac12$ for every $q\le2^{64}$. The standard expander-walk Chernoff and hitting
bounds \cite{hlw2006} therefore apply to walks on $\Gq$ driven by these samples. A walk of length $t$ from a
uniform start uses $\log_2\varphi(q)+t\,s$ random bits instead of $(t+1)\log_2\varphi(q)$. Drawing the uniform
start does require sampling $\Gq$, e.g.\ by rejection with gcd tests.

\paragraph{What this is not.}
These are statistical guarantees against multiplicative characters of $\Gq$. They give no
indistinguishability against other tests (e.g.\ additive structure of $n$, or $n$ itself), and no
cryptographic security.
